\documentclass[acmtog]{acmart}
\usepackage{graphicx}
\usepackage{natbib}
\usepackage{listings}
\usepackage{bm}
\usepackage{amsmath}
\usepackage{color}
\usepackage{underscore}
\usepackage{xcolor}
\usepackage{soul}
\usepackage{subcaption}

\definecolor{blve}{rgb}{0.3372549 , 0.61176471, 0.83921569}
\definecolor{gr33n}{rgb}{0.29019608, 0.7372549, 0.64705882}
\makeatletter
\lst@InstallKeywords k{class}{classstyle}\slshape{classstyle}{}ld
\makeatother
% 关键优化：添加自动换行、字符边界换行、禁用行号溢出
\lstset{
	language=C++,
	basicstyle=\small\ttfamily,  % 缩小字体，增加每行可显示字符数
	keywordstyle=\color{blve}\ttfamily,
	stringstyle=\color{red}\ttfamily,
	commentstyle=\color{magenta}\ttfamily,
	morecomment=[l][\color{magenta}]{\#},
	classstyle = \bfseries\color{gr33n}, 
	tabsize=2,
	breaklines=true,  % 启用自动换行
	breakatwhitespace=false,  % 允许在非空白字符处换行（避免长单词/变量名溢出）
	breakindent=2em,  % 换行后缩进2字符，保持可读性
	columns=flexible,  % 灵活列宽，适配换行
	showstringspaces=false,  % 不显示字符串中的空格标记
	frame=tb,  % 上下边框，优化显示
	aboveskip=1mm,  % 代码块上间距
	belowskip=1mm   % 代码块下间距
}

% Title portion
\title{Assignment 1:\\ {Exploring OpenGL Programming}} 

\author{Name: \quad Yixuan Yan\\ student number:2023533152\
	\\email: \quad \texttt{yanyx2023@shanghaitech.edu.cn}}

% Document starts
\begin{document}
	\maketitle
	
	\vspace*{2 ex}
	
	\section{Introduction}
	This assignment focuses on implementing core components of a ray-tracing renderer, including ray-triangle intersection tests, BVH (Bounding Volume Hierarchy) acceleration structures, direct illumination integrators, refractive materials, and anti-aliasing. The final goal is to render a Cornell Box scene with physically plausible lighting and refraction effects.
	
	Completed tasks include:
	\begin{itemize}
		\item Implement ray-triangle and ray-AABB intersection tests;
		\item Construct a BVH acceleration structure to optimize ray-tracing efficiency;
		\item Develop a direct illumination integrator supporting refractive materials;
		\item Integrate shadow testing and multi-sample anti-aliasing;
		\item Validate the implementation by rendering the Cornell Box scene.
	\end{itemize}
	
	
	\section{Implementation Details}
	\subsection{Task 1: Geometric Intersection Tests}
	Implement ray-triangle and ray-AABB (Axis-Aligned Bounding Box) intersection tests, with relevant code located in \texttt{src/accel.cpp}.
	
	\subsubsection{Ray-Triangle Intersection (\hl{\textit{TriangleIntersect}})}
	Implemented using the geometric method:
	\begin{lstlisting}
		bool TriangleIntersect(const Ray& ray, 
		const Vec3f& v0, const Vec3f& v1,
		const Vec3f& v2,
		Float& t, Float& u, Float& v) {
			Vec3f edge1 = v1 - v0;
			Vec3f edge2 = v2 - v0;
			Vec3f h = Cross(ray.d, edge2);
			Float a = Dot(edge1, h);
			if (a > -1e-8 && a < 1e-8) return false; 
			// Ray is parallel to the triangle plane
			Float f = 1.0f / a;
			Vec3f s = ray.o - v0;
			u = f * Dot(s, h);
			if (u < 0.0 || u > 1.0) return false;
			Vec3f q = Cross(s, edge1);
			v = f * Dot(ray.d, q);
			if (v < 0.0 || u + v > 1.0) return false;
			t = f * Dot(edge2, q);
			return t > 1e-8; // Exclude reverse rays
		}
	\end{lstlisting}
	Core logic: Calculate the intersection point of the ray with the triangle plane using cross products, then verify if the point lies within the triangle using barycentric coordinates.
	
	\subsubsection{Ray-AABB Intersection (\hl{\textit{AABB::intersect}})}
	Implemented using the Slab method:
	\begin{lstlisting}
		bool AABB::intersect(const Ray& ray, 
		Float& t_min, Float& t_max) const {
			for (int i = 0; i < 3; ++i) {
				Float invD = 1.0f / ray.d[i];
				Float t0 = (min[i] - ray.o[i]) * invD;
				Float t1 = (max[i] - ray.o[i]) * invD;
				if (invD < 0.0f) std::swap(t0, t1);
				t_min = std::max(t0, t_min);
				t_max = std::min(t1, t_max);
				if (t_min > t_max) return false;
			}
			return t_max > 1e-8;
		}
	\end{lstlisting}
	Core logic: Compute intersection points of the ray with the AABB along the three axes, then determine if a valid intersection interval exists by comparing the maximum near intersection and minimum far intersection.
	
	\subsection{Task 2: BVH Construction and Traversal}
	Implement BVH tree construction (\hl{\textit{BVHTree<_>::build}} in \texttt{bvh\_tree.h}) using a recursive partitioning strategy:
	\begin{lstlisting}
		template <typename Primitive>
		void BVHTree<Primitive>::build(
		std::vector<Primitive>& primitives) {
			root = std::make_unique<BVHNode>();
			root->primitives = std::move(primitives);
			// Compute node bounding box
			root->bounds = ComputeBounds(root->primitives);
			if (root->primitives.size() <= max_leaf_size) 
			return;
			
			// Select partitioning axis
			int axis = root->bounds.maxExtent();
			// Sort primitives along the axis
			std::sort(root->primitives.begin(), 
			root->primitives.end(),
			[axis](const Primitive& a, const Primitive& b) {
				return a.getBounds().centroid()[axis] < 
				b.getBounds().centroid()[axis];
			});
			// Partition primitives
			auto mid = root->primitives.begin() + 
			root->primitives.size() / 2;
			std::vector<Primitive> left_primitives(
			root->primitives.begin(), mid);
			std::vector<Primitive> right_primitives(
			mid, root->primitives.end());
			root->primitives.clear(); // Internal nodes do not store primitives
			
			// Recursively build subtrees
			root->left = std::make_unique<BVHNode>();
			buildRecursive(root->left, std::move(left_primitives));
			root->right = std::make_unique<BVHNode>();
			buildRecursive(root->right, std::move(right_primitives));
		}
	\end{lstlisting}
	Core logic:
	1. Compute the bounding box for the node;
	2. If the number of primitives exceeds the threshold, partition along the longest axis of the bounding box;
	3. Recursively build left and right subtrees to form the BVH hierarchy.
	
	\subsection{Task 3: Direct Illumination Integrator with Shadow Testing}
	Implement the \hl{\textit{IntersectionTestIntegrator}} (in \texttt{src/integrator.cpp}) to handle ray tracing, direct illumination calculation, and shadow testing.
	
	\subsubsection{Multi-Sample Anti-Aliasing (\hl{\textit{IntersectionTestIntegrator::render}})}
	Generate multiple samples per pixel to reduce aliasing:
	\begin{lstlisting}
		for (int dx = 0; dx < resolution.x; dx++) {
			for (int dy = 0; dy < resolution.y; dy++) {
				for (int sample = 0; sample < spp; sample++) {
					Vec2f pixel_sample = sampler.getPixelSample(); 
					// Sub-pixel offset
					auto ray = camera->generateDifferentialRay(
					pixel_sample.x, pixel_sample.y);
					Vec3f L = Li(scene, ray, sampler);
					camera->getFilm()->commitSample(
					pixel_sample, L);
				}
			}
		}
	\end{lstlisting}
	
	\subsubsection{Ray Tracing and Material Handling (\hl{\textit{IntersectionTestIntegrator::Li}})}
	Trace rays and handle refractive materials:
	\begin{lstlisting}
		Vec3f IntersectionTestIntegrator::Li(
		ref<Scene> scene, DifferentialRay& ray, 
		Sampler& sampler) const {
			Vec3f color(0.0);
			DifferentialRay current_ray = ray;
			for (int i = 0; i < max_depth; ++i) {
				SurfaceInteraction isect;
				if (!scene->intersect(current_ray, isect)) break;
				
				isect.wo = -current_ray.direction;
				auto* bsdf = isect.bsdf;
				if (auto* refract = dynamic_cast<const PerfectRefraction*>(bsdf)) {
					// Handle refraction
					Float pdf;
					Vec3f wi = refract->sample(isect, sampler, &pdf);
					Vec3f spawn_origin = isect.p + wi * 1e-3f; 
					// Avoid self-intersection
					current_ray = DifferentialRay(spawn_origin, wi);
					continue;
				}
				
				if (dynamic_cast<const IdealDiffusion*>(bsdf)) {
					// Compute direct illumination
					color = directLighting(scene, isect);
					break;
				}
				break;
			}
			return color;
		}
	\end{lstlisting}
	
	\subsubsection{Direct Illumination and Shadow Testing (\hl{\textit{IntersectionTestIntegrator::directLighting}})}
	Calculate diffuse illumination and check for shadows:
	\begin{lstlisting}
		Vec3f IntersectionTestIntegrator::directLighting(
		ref<Scene> scene, SurfaceInteraction& isect) const {
			Vec3f light_dir = Normalize(
			point_light_position - isect.p);
			Float dist_to_light = Norm(
			point_light_position - isect.p);
			// Generate shadow ray
			Vec3f shadow_origin = isect.p + isect.normal * 1e-3f;
			DifferentialRay shadow_ray(
			shadow_origin, light_dir, 1e-3f, 
			dist_to_light - 1e-3f);
			SurfaceInteraction shadow_isect;
			if (scene->intersect(shadow_ray, shadow_isect)) 
			return Vec3f(0.0); // Occluded
			
			// Lambertian diffuse calculation
			auto* diffuse_bsdf = dynamic_cast<const IdealDiffusion*>(isect.bsdf);
			Vec3f albedo = diffuse_bsdf->evaluate(isect);
			Float cos_theta = std::max(
			Dot(light_dir, isect.normal), 0.0f);
			Float attenuation = 1.0f / (dist_to_light * dist_to_light); 
			// Inverse-square law
			return albedo * cos_theta * point_light_flux * 
			attenuation * 0.1f;
		}
	\end{lstlisting}
	
	\subsection{Task 4: Perfect Refraction Material }
	Implement Snell's Law to compute refraction directions (in \texttt{src/bsdf.cpp}):
	\begin{lstlisting}
		Vec3f PerfectRefraction::sample(
		SurfaceInteraction& isect, Sampler& sampler, 
		Float* pdf) const {
			Vec3f n = isect.shading.n;
			Float cosi = Dot(isect.wo, n);
			bool entering = cosi > 0;
			Float eta_i = entering ? 1.0f : eta;
			Float eta_t = entering ? eta : 1.0f;
			if (!entering) n = -n;
			
			Vec3f refracted_dir;
			if (!Refract(isect.wo, n, eta_i / eta_t, refracted_dir)) {
				refracted_dir = Reflect(isect.wo, n); 
				// Total internal reflection
			}
			isect.wi = refracted_dir;
			if (pdf) *pdf = 1.0f;
			return Vec3f(1.0f);
		}
	\end{lstlisting}
	
	\section{Results}
	\begin{figure}[htbp]
		\centering
		\begin{subfigure}[b]{0.4\textwidth}  
			\centering
			\includegraphics[width=\textwidth]{cbox_no_light.png} 
			\caption{cbox_no_light}
			\label{fig:sub1}
		\end{subfigure}
		\hfill
		\begin{subfigure}[b]{0.4`\textwidth}
			\centering
			\includegraphics[width=\textwidth]{cbox_no_light_refract.png}
			\caption{cbox_no_light_refract}
			\label{fig:sub2}
		\end{subfigure}
		
	\end{figure}
	
	
\end{document}